\begin{abstract}
We determine the minimum number of proper Euclidean circles containing at least three selected points in a nondegenerate set of $n$ planar points, where the points are distinct, not all collinear, and not all concyclic.  This is the Elliott--Purdy--Smith interpretation of Erd\H{o}s Problem 506; the literal wording without the collinearity exclusion is degenerate.  We prove $f(8)=17$ and
\[
 f(n)=1+\binom{n-1}{2}-\left\lfloor\frac{n-1}{2}\right\rfloor
 \qquad(n\ge 9).
\]
The previously unresolved finite cases are $f(12)=51$ and $f(13)=61$.  The proof combines generalized-circle block partitions, inversion, line-arrangement inequalities, exact rational certificates, generalized allowable sequences, exhaustive block-packing searches, and a projective nonrealizability argument.  Every finite computation is supplied in a self-contained verification archive with source code, finite inputs, manifests, retained exact outputs, and clean replay commands.  No network access or proprietary optimizer is required for the canonical replay.
\end{abstract}

\subjclass[2020]{Primary 52C10, 52C30; Secondary 05B35, 68V20}
\keywords{Erd\H{o}s problem, discrete geometry, incidence geometry, circles determined by points, computer-assisted proof, pseudoline arrangement, exact enumeration}

\maketitle

\begin{center}
\fcolorbox{black}{gray!8}{\parbox{0.91\linewidth}{\small
\textbf{Preprint status.} This manuscript reports a computer-assisted proof and has not yet been journal-refereed.  Independent mathematical review and clean-room reimplementation of the principal finite searches are explicitly invited.}}
\end{center}

